\documentclass[11pt,a4paper]{article}
\usepackage[utf8]{inputenc}
\usepackage[T1]{fontenc}
\usepackage[turkish]{babel}
\usepackage{lmodern}
\usepackage{microtype}
\usepackage{amsmath, amssymb, amsfonts}
\usepackage[margin=2.4cm]{geometry}
\usepackage{setspace}
\usepackage{url}
\usepackage{xurl}
\usepackage{booktabs}
\usepackage{tabularx}
\usepackage{graphicx}
\onehalfspacing

\title{\textbf{Stationary Strings Mektup 1: Üç-Foton Pozitronyum Spektrumu --- arXiv:2605.21289 Denk.~(12)'nin Düzeltmesi ve Bir Saatlik Laboratuvar Testi}}
\author{Fahri Meral}
\date{20 Ekim 2026}

\begin{document}
\shorthandoff{=}

\maketitle

\begin{abstract}
arXiv:2605.21289'da (Denk.~(12)) kullanılan orto-pozitronyum üç-foton
spektral şekli, 1949 Ore--Powell ifadesinin logaritmasız (ln'siz) hâlidir
ve fiziksel olarak tutarsızdır: normalize değildir
($\int_1^{511}\Omega\,\mathrm{d}E = 2343.47$, $1$ olmalıdır) ve düşük
enerji davranışı $E^{-2}$ olarak ıraksar; yani yok olması gereken
foton sayısından fazlasını öngörür. Doğru ifade, 1949 makalesinin
$\ln(1-x)$ çarpanlarını taşır, tam olarak normaldir
($\int_0^{511}\Omega\,\mathrm{d}E = 1$, $5\times10^{-13}$ hassasiyetle)
ve $x = E/m_e \to 0$ iken $\Omega(x)\sim\tfrac{5}{3}x$ ile sönümlenir.
$5\,\mu$Ci $^{22}$Na kaynağı, silika aerojel ve küçük bir HPGe dedektörle
yapılacak \emph{tek saatlik} bir ölçümün iki formu $7\times10^{5}$'den
büyük standart sapma ile ayırdığını gösteriyoruz; yani düzeltme doğrudan
bir laboratuvar tezgâhında sınanabilir. Burada raporlanan ölçüm
\emph{yapılmamıştır, önerilmiştir}: aşağıdaki sayılar protokolün öngörülen
sonucudur ve düzeltmenin kendisi için kanıt, yayımlanmış denklemin 1949
aslına karşı kaynak metni denetimidir. Protokolün tamamı, üç aday
spektrum ve tüm sayısal makine bu mektupla birlikte yayımlanmaktadır.
\end{abstract}

\section{Sorun}
Ortogonal-pozitronyumun (o-Ps) üç-foton spektral şekli
$\Omega(E_{\rm em}/m_e)$, arXiv:2605.21289'daki~\cite{sy26} ekstragalaktik
511~keV çizgisi modelinde spektral şekil üzerinden girer. Yayımlanan
Denk.~(12) şu ifadedir:
\begin{equation}
\label{eq:pub}
\Omega(E) = \frac{2}{m_e(\pi^2-9)}\left[
\frac{E(m_e-E)}{(2m_e-E)^2} + \frac{2m_e-E}{E}
+ \frac{2m_e(m_e-E)}{E^2} - \frac{2m_e(m_e-E)^2}{(2m_e-E)^3}\right],
\end{equation}
$m_e = 511$~keV ile. Bu, 1949 Ore--Powell ifadesinin ln'siz
varyantıdır~\cite{ore49}: aslında taşıdığı $\ln(1-x)$ çarpanları
yoktur. İki bağımsız belirti doğar ve ikisi de nicel kullanım için
ölümcüldür:
\begin{enumerate}
    \item \emph{Normalizasyon.} Denk.~\eqref{eq:pub} yerine konduğunda
    $\int_{1}^{511}\Omega\,\mathrm{d}E = 2343.47$ çıkar, $1$ değil.
    İntegrali kendi normalizasyonunun üç büyüklük mertebesi kadar
    aşan bir spektral şekil olasılık yoğunluğu olamaz.
    \item \emph{Düşük-enerji ıraksaması.} Üçüncü terim $2m_e^2/E^2$
    gibi davranır: şekil $E\to 0$ iken karesel ıraksar, oysa üç-foton
    faz uzayı orada sönümlenmek zorundadır. Doğru ifade ise
    $\Omega(x)\sim\tfrac{5}{3}x$, yani $O(E)\propto E$ sağlar.
\end{enumerate}
Fiziksel olarak doğru form, üçüncü ve dördüncü terimleri
\begin{equation}
\label{eq:corr}
\frac{2m_e(m_e-E)}{E^2} - \frac{2m_e(m_e-E)^2}{(2m_e-E)^3}
\ \longrightarrow\
\frac{2m_e(m_e-E)}{(2m_e-E)^2} - \frac{2m_e(m_e-E)^2}{(2m_e-E)^3}
\ +\ \ln(1-x)\ \text{çarpanları},
\end{equation}
şeklinde düzeltir; yani 1949 ifadesinin $\ln(1-x)$ çarpanları geri
konmuş hâlidir: tam olarak normaldir
($\int_0^{511}\Omega\,\mathrm{d}E = 1$, $5\times10^{-13}$ hassasiyetle)
ve $E\to 0$'da sönümlenir. Kısmi bir onarım --- yalnız üçüncü
terimde $E^2 \to (2m_e-E)^2$ ikamesi --- ikinci terimin $1/E$ kutbunu
bırakır ve $\int_{1}^{511}\Omega\,\mathrm{d}E = 27.147$ verir; o da $1$
değildir.

\section{Bu neden tek bir denklemden ibaret değil}
511~keV çizgi şekli, difüz pozitron annihilation akısını haritalamak ve
karanlık madde annihilation kanallarına kısıtlar koymak için
kullanılır~\cite{sy26,yoneda25}. Yanlış bir spektral şekil, doğrudan
çıkarılan akıya ve buradan türetilen karanlık madde sınırlarına birim
mertebesinde hata olarak yayılır. Hata kozmetik değildir: yayımlanmış bir
astrofizik çıkarımda birim mertebesinde hatadır ve normalizasyon
kontrol edilmedikçe görünmez.

\section{Bir saatlik laboratuvar testi}
Üç aday şekil bir bant-oranı ölçümüyle ayrılır. $5\,\mu$Ci
($1.85\times10^5$ Bq) $^{22}$Na kaynağı, silika aerojel hedef
($f_{\rm Ps} = 0.55$), katı açı $\times$ verim $\epsilon = 0.02$ ve
$3600$~s toplama süresi, $\sigma = 0.85$~keV'lik bir HPGe ile
(511~keV'de FWHM $\approx 2.00$~keV) alınır. Dedekte edilen üç-foton
sayısı
$N_{3\gamma} = A f_{\beta^+} f_{\rm Ps} \tfrac{3}{4}\epsilon\,t = 4.96\times 10^{6}$
ve ayırt edici büyüklük, spektral şeklin kendisine ait olan bant oranıdır:
$R = N_{[10,100]}/N_{[350,490]}$.

\begin{table}[!htbp]
\centering\footnotesize
\setlength{\tabcolsep}{4pt}
\caption{\textbf{Üç aday spektrum için bir saatlik ölçümün öngörülen sonucu.}
Ölçüm yapılmamıştır; bunlar protokolün öngörüleridir. Yalnız Poisson
hataları.}
\label{tab:na22}
\vspace{0.1cm}
\begin{tabular}{@{}lrrr@{}}
\toprule
Form & $N_{[10,100]}$ & $N_{[350,490]}$ & $R$ \\
\midrule
Düzeltilmiş 1949 (exact) & $1.847\times 10^{5}$ & $2.195\times 10^{6}$ & $0.08414\pm 0.00020$ \\
Yayımlanan (ln'siz) & $1.047\times 10^{9}$ & $6.570\times 10^{6}$ & $159.37\pm 0.06$ \\
Kısmi onarım & $5.112\times 10^{7}$ & $5.494\times 10^{6}$ & $9.305\pm 0.004$ \\
\bottomrule
\end{tabular}
\end{table}

Ayrım kesindir: düzeltilmiş ve yayımlanan formlar $R$'de
$7.8\times 10^{5}$ standart sapma ile ayrılır; ayrıca yayımlanan form
düşük-enerji bandında düzeltilmiş formdan $5.67\times 10^{3}$ kat fazla
sayım öngörür --- tek başına $2.4\times10^{6}\,\sigma$'luk bir fazlalık ve
gerçekçi hiçbir dedektörü doyuracak bir sayım hızı. Bir dakikalık
toplama bile $10^{5}\,\sigma$ mertebesinde ayrımı korur. Banddaki
baskın fiziksel arka plan sınanan üç-foton kontinuumunun kendisidir;
bu yüzden ayırt edici büyüklük alışılmış anlamda arka planla
sınırlı değildir. Artık arka planda on kat belirsizlik olsa bile ayrım
$10^{4}\,\sigma$'nun üzerindedir.

\section{Ne iddia ediliyor, ne edilmiyor}
Düzeltmenin kanıtı bir \emph{kaynak denetimidir}: yayımlanmış denklem
terim terim 1949 aslıyla ve literatürdeki bağımsız kontrollerle
karşılaştırılmıştır. Burada önerilen laboratuvar testi bu denetimin
çürütme testidir ve \emph{yapılmamıştır}; yürütülebilmesi ve
karşılaştırılabilmesi için öngörülen sonucu raporluyoruz. Bu mektubun
hiçbir iddiası, serinin modeline bağlı değildir; bu yüzden seriden ayrı
olarak sunulmaktadır.

\section{Yeniden-üretilebilirlik}
Tüm makine bu mektupla birlikte yayımlanmıştır: üç spektral şekil, bant
integralleri, dedekte edilen sayım ve ayrım hesabı ve bağımsız sayısal
kontroller (exact formun integrali $5\times10^{-13}$ ile $1$'e eşittir ve
düşük-enerji davranışının $E$'ye doğrusal olduğu doğrulanır).
Çalıştırma talimatları ek depoların README dosyasındadır.

\begin{thebibliography}{99}
    \bibitem{ore49} Ore, A. \& Powell, J.~L. (1949). Three-photon
    annihilation of an electron-positron pair. \textit{Phys.\ Rev.\ \textbf{75}},
    1696--1699.
    \bibitem{sy26} Siegert, T. \& Yoneda, H. (2026). On the possibility of
    an extragalactic positron annihilation signal.
    arXiv:2605.21289.
    \bibitem{yoneda25} Yoneda, H., Siegert, T. \& Mittal, S. (2025).
    Imaging the positron annihilation line with 20 years of
    INTEGRAL/SPI observations. \textit{Astron.\ Astrophys.\ \textbf{702}},
    A220. arXiv:2509.01066.
    \bibitem{ps95} Positronium Working Group (1995). Report of the
    IUPAP positronium panel. \textit{Prog.\ Part.\ Nucl.\ Phys.\ \textbf{17}},
    219--287.
\end{thebibliography}

\end{document}
